% Compact body table: key coefficients discussed in the text.
% Full list: Table~\ref{tab:app-secsize-recovery} in the Online Appendix.
\begin{table}[H]
\centering
\caption{Selected section and size coefficients in the linear recovery comparison.}\label{tab:secsize-recovery}
\footnotesize
\begin{tabular}{@{}llp{3.6cm}rr@{}}
\toprule
Segment & Code & Term & Coef. & HC3 95\% CI \\
\midrule
Declared & C & Manufacturing & 1.293 & 0.485 to 2.100 \\
Declared & H & Transportation and storage & 1.258 & 0.587 to 1.929 \\
Declared & J & Information and communication & 0.088 & $-0.484$ to 0.660 \\
Declared & 03 & Small & 0.187 & $-0.345$ to 0.720 \\
Declared & 05 & Other sizes & 0.627 & 0.095 to 1.159 \\
Assessed & C & Manufacturing & $-0.009$ & $-0.039$ to 0.021 \\
Assessed & H & Transportation and storage & $-0.007$ & $-0.045$ to 0.032 \\
Assessed & J & Information and communication & 0.012 & $-0.032$ to 0.056 \\
Assessed & 03 & Small & 0.036 & 0.007 to 0.066 \\
Assessed & 05 & Other sizes & 0.069 & 0.038 to 0.101 \\
\bottomrule
\end{tabular}

\vspace{0.4em}
\begin{minipage}{0.95\textwidth}
\footnotesize
Notes: Ordinary least squares of payments divided by stock, inside each segment. This equation is a comparison with the payment logit. It is not the stock equation. Declared uses 194,709 certificates. Assessed uses 5,092 certificates. Trade (section G) and micro (size code 01) are the reference levels and have no coefficient. Coefficients and intervals are rounded to three decimal places. The full coefficient list is Table~\ref{tab:app-secsize-recovery}.
\end{minipage}
\end{table}
